\documentclass[12pt,a4paper]{article}

% ---- Encodage et langue ----
\usepackage[utf8]{inputenc}
\usepackage[T1]{fontenc}
\usepackage[french]{babel}

% ---- Mise en page ----
\usepackage[top=2.5cm, bottom=2.5cm, left=2.5cm, right=2.5cm]{geometry}
\usepackage{setspace}
\onehalfspacing

% ---- Typographie ----
\usepackage{lmodern}
\usepackage{microtype}

% ---- Mathématiques ----
\usepackage{amsmath, amssymb, amsthm}

% ---- Algorithmique ----
\usepackage[ruled,vlined,linesnumbered,french]{algorithm2e}

% ---- Graphiques et tableaux ----
\usepackage{graphicx}
\usepackage{booktabs}
\usepackage{array}
\usepackage{xcolor}
\usepackage{colortbl}

% ---- Liens et méta-données ----
\usepackage[colorlinks=true, linkcolor=blue!70!black, urlcolor=blue!70!black,
            citecolor=green!50!black]{hyperref}

% ---- Code source ----
\usepackage{listings}
\lstset{
  basicstyle=\ttfamily\small,
  breaklines=true,
  frame=single,
  backgroundcolor=\color{gray!10},
  keywordstyle=\color{blue},
  commentstyle=\color{green!50!black},
  stringstyle=\color{red!70!black},
  numbers=left,
  numberstyle=\tiny\color{gray},
  language=Python
}

% ---- Références ----
\usepackage{cite}

% =========================================================
\begin{document}

% =========================================================
%  PAGE DE GARDE
% =========================================================
\begin{titlepage}
  \centering
  \vspace*{1cm}

  {\Huge\bfseries GraphBench Challenge\par}
  \vspace{0.4cm}
  {\Large Réfutation automatique de conjectures en théorie des graphes\par}

  \vspace{1.5cm}
  \rule{\linewidth}{0.4pt}
  \vspace{0.8cm}

  {\large\bfseries Étudiants :\par}
  \vspace{0.3cm}
  {\large Martim MOREIRA ALVES\par}
  {\large Noé ARNAUD\par}

  \vspace{0.8cm}
  {\large\bfseries Formation :\par}
  {\large Master 1 MIAGE — Université Paris Dauphine\par}

  \vspace{0.8cm}
  {\large\bfseries Date :\par}
  {\large 12 mai 2026\par}

  \vspace{0.8cm}
  {\large\bfseries Dépôt GitHub :\par}
  \href{https://github.com/AlvesMartim/OC_Graph}{\texttt{https://github.com/AlvesMartim/OC\_Graph}}

  \vspace{1.2cm}
  \rule{\linewidth}{0.4pt}
  \vspace{0.8cm}

  \begin{center}
    \begin{tabular}{cc}
      \toprule
      \textbf{Conjectures réfutées} & \textbf{Score final (secondes)} \\
      \midrule
      \textbf{99 / 100} & \textbf{221,51 s} \\
      \bottomrule
    \end{tabular}
  \end{center}

  \vfill
  {\small Rapport de projet — M1 MIAGE}
\end{titlepage}

% =========================================================
%  TABLE DES MATIÈRES
% =========================================================
\tableofcontents
\newpage

% =========================================================
%  1. INTRODUCTION ET OBJECTIF
% =========================================================
\section{Introduction et objectif}

L'objectif de ce projet est simple : on nous donne 100 conjectures sur les graphes,
et on doit en réfuter le plus possible automatiquement.  Une conjecture prend la forme
d'une inégalité entre deux mesures d'un graphe (appelées \emph{invariants}) :

\[
  Y(G) \;\leq\; f\!\bigl(X(G)\bigr)
  \qquad \text{ou} \qquad
  Y(G) \;\geq\; f\!\bigl(X(G)\bigr)
\]

où $X$ et $Y$ sont par exemple le degré maximum, le nombre de triangles, ou la
proximité, et $f$ est un polynôme donné.  Réfuter une conjecture, c'est trouver
un graphe $G$ pour lequel l'inégalité est fausse.

Pour mesurer à quel point un graphe ``viole'' une conjecture, on calcule :

\[
  \mathrm{violation}(G) =
  \begin{cases}
    Y(G) - f\!\bigl(X(G)\bigr) & \text{si signe } \leq \\
    f\!\bigl(X(G)\bigr) - Y(G) & \text{si signe } \geq
  \end{cases}
\]

Dès que ce score est positif, on a trouvé un contre-exemple.  Le score final du
projet est la \textbf{somme des temps de réfutation}, avec une pénalité de $120$~s
pour chaque conjecture qu'on n'arrive pas à réfuter.

Nous avons développé deux approches.  La \textbf{Partie~1} est une heuristique
de recherche locale : on part de graphes connus ou aléatoires, et on les modifie
progressivement pour augmenter le score de violation.  La \textbf{Partie~2}
utilise \textbf{FunSearch} : on demande à un LLM de générer des fonctions de
guidage plus intelligentes, qu'on teste et améliore par itérations.

% =========================================================
%  2. HEURISTIQUE SIMPLE
% =========================================================
\section{Heuristique simple}

\subsection{Vue d'ensemble}

Pour chaque conjecture, l'heuristique suit trois étapes dans l'ordre
(Figure~\ref{fig:pipeline}).  L'idée générale est de d'abord tester rapidement
des graphes ``classiques'', puis si ça ne suffit pas, de lancer une recherche
plus longue en modifiant des graphes au fil du temps.

\begin{figure}[h]
  \centering
  \begin{tabular}{|c|c|c|}
    \hline
    \textbf{Phase 1} & \textbf{Phase 2} & \textbf{Phase 3} \\
    Passe atlas ($\leq 5$~s) & Population initiale & Boucle locale (60~s) \\
    \hline
    Graphes $\leq 7$ nœuds & Seeds structurées & Tournoi $k=3$ \\
    Arbres 8--12 nœuds & Graphes aléatoires & Mutation ciblée 70\% \\
    Familles paramétriques & (10 graphes) & Liste taboue + kick \\
    \hline
  \end{tabular}
  \caption{Pipeline de l'heuristique simple.}
  \label{fig:pipeline}
\end{figure}

\subsection{Phase 1 : passe atlas}

Avant de faire quoi que ce soit de complexe, on teste une grande liste de graphes
connus pour voir si l'un d'eux est déjà un contre-exemple :

\begin{itemize}
  \item \textbf{Atlas NetworkX} : tous les graphes connexes à au plus 7 sommets
        (1\,252 graphes), testés en quelques millisecondes.
  \item \textbf{Arbres} de 8 à 12 sommets, si la conjecture porte sur les arbres.
  \item \textbf{Familles classiques} pour $n$ de 8 à 25 : chemins, cycles, étoiles,
        graphes complets, roues, chenilles, doubles-étoiles, bipartis, $K_n +$ bras.
\end{itemize}

Cette étape résout à elle seule \textbf{72 des 99 conjectures réfutées}.
Beaucoup de conjectures ont des contre-exemples très simples, comme $K_2$
(le graphe à 2 sommets reliés par une arête), trouvé en moins de 0,1~s.

\subsection{Phase 2 : population initiale et seeds structurées}

Si l'atlas n'a rien trouvé, on construit une population de départ de 10 graphes
pour amorcer la recherche.  Une partie de ces graphes est choisie \emph{intelligemment}
selon les invariants de la conjecture :

\begin{itemize}
  \item Si la conjecture parle de \texttt{remoteness} : on part de graphes
        $K_n +$ bras (une clique avec un chemin accroché), qui ont une remoteness
        très basse.
  \item Si la conjecture parle de \texttt{proximity} : on part de doubles-étoiles
        très déséquilibrées ou de ``brooms'' (chemin + étoile à un bout), qui
        étirent le graphe et font baisser la proximity.
\end{itemize}

Le reste est généré aléatoirement en respectant la classe de graphes demandée
(\texttt{connected}, \texttt{tree}, ou \texttt{claw\_free}).

\subsection{Phase 3 : recherche locale avec diversification}

La boucle principale modifie les graphes de la population un par un,
en gardant ceux qui progressent vers une violation positive.
Trois mécanismes évitent de rester bloqué au même endroit :

\begin{algorithm}[H]
\SetAlgoLined
\KwIn{Conjecture $C$, budget temps $T = 60$~s}
\KwOut{Contre-exemple $G^*$ ou échec}
\BlankLine
$\text{tabu} \leftarrow \emptyset$\;
$\text{pop} \leftarrow \text{GenerateInitialGraphs}(C)$\;
Évaluer chaque $G \in \text{pop}$, conserver les 10 meilleurs\;
\BlankLine
\While{temps écoulé $< T$}{
  $G \leftarrow \text{TournoiSelection}(\text{pop},\, k=3)$\;
  \eIf{stagnation $>$ KICK (100 iter)}{
    $H \leftarrow \text{KickMutate}(G,\, 4\text{ mutations})$\;
  }{
    $H \leftarrow \text{TargetedMutate}(G,\, C)$\;
  }
  $H \leftarrow \text{Repair}(H,\, \text{classe}(C))$\;
  \If{$\mathrm{graph6}(H) \in \text{tabu}$}{\textbf{continue}\;}
  Ajouter $\mathrm{graph6}(H)$ dans tabu\;
  Calculer $\mathrm{violation}(H, C)$\;
  \If{$\mathrm{violation}(H) > 0$}{\Return $H$\;}
  Mettre à jour \text{pop}\;
  \If{stagnation $>$ RESTART (500 iter)}{
    $\text{pop} \leftarrow \text{GenerateInitialGraphs}(C)$\;
    stagnation $\leftarrow 0$\;
  }
}
\Return échec\;
\caption{Recherche locale avec liste taboue, kick et redémarrage}
\end{algorithm}

\paragraph{Sélection par tournoi.}
On tire 3 graphes au hasard dans la population et on choisit le meilleur.
Cela favorise les bons graphes sans pour autant ignorer les autres.

\paragraph{Mutations ciblées.}
70\% du temps, la modification est \emph{guidée} par la conjecture : si on veut
augmenter $Y$, on applique une mutation conçue pour ça.  Par exemple :
\begin{itemize}
  \item Augmenter \texttt{triangle\_number} : on cherche deux voisins d'un
        même sommet qui ne sont pas encore reliés, et on ajoute l'arête.
  \item Augmenter \texttt{diameter} : on accroche un nouveau sommet à une
        feuille existante pour allonger le graphe.
  \item Diminuer \texttt{proximity} : on ajoute un sommet à l'endroit le plus
        ``loin'' du reste du graphe, pour l'étirer encore plus.
  \item Augmenter \texttt{clique\_number} : on connecte un sommet extérieur
        à tous les membres de la plus grande clique actuelle.
\end{itemize}
Les 30\% restants sont des mutations aléatoires (ajout/suppression d'arête ou de sommet)
pour maintenir la diversité.

\paragraph{Liste taboue.}
Chaque graphe visité est converti en une chaîne \texttt{graph6} unique et mémorisé.
Si on retombe sur un graphe déjà vu, on l'ignore.  Cela évite de tourner en rond.

\paragraph{Kick et redémarrage.}
Si le score ne progresse pas depuis 100 itérations, on applique 4 mutations
d'un coup pour sortir du blocage (\emph{kick}).  Si ça ne suffit toujours pas
après 500 itérations, on repart de zéro avec une nouvelle population.

\subsection{Calcul des invariants}

On calcule 25 invariants différents, mais seulement ceux dont la conjecture
a besoin (calcul \emph{lazy}).  Les deux invariants de distance les plus utilisés
sont définis ainsi : si $T_v$ est la somme des distances de $v$ à tous les autres
sommets,

\[
  \mathrm{proximity}(G) = \frac{n-1}{\max_{v} T_v}, \qquad
  \mathrm{remoteness}(G) = \frac{n-1}{\min_{v} T_v}.
\]

Les invariants spectraux (valeur propre maximale, connectivité algébrique) sont
calculés avec NumPy.  Les invariants combinatoires comme le nombre de domination
utilisent des algorithmes de NetworkX.

\subsection{Réparation des graphes}

Une mutation peut casser une propriété requise (par exemple déconnecter le graphe).
Après chaque mutation, on répare automatiquement :

\begin{itemize}
  \item \texttt{connected} : on relie les composantes déconnectées par des arêtes aléatoires.
  \item \texttt{tree} : on calcule un arbre couvrant (algorithme de Kruskal).
  \item \texttt{claw\_free} : on supprime les arêtes qui créent des griffes $K_{1,3}$.
\end{itemize}

% =========================================================
%  3. ARCHITECTURE FUNSEARCH
% =========================================================
\section{Architecture FunSearch}

\subsection{Principe général}

FunSearch~\cite{romera2024mathematical} est une technique proposée par Google
DeepMind où un LLM génère du code, ce code est testé automatiquement, et les
meilleurs résultats sont réinjectés dans le prompt pour guider la prochaine
génération.  Ici, on l'applique de la façon suivante : au lieu de coder à la
main comment guider la recherche, on demande au LLM de générer une
\emph{fonction de score} Python qui remplace notre score de violation dans la
boucle de recherche.

\subsection{Boucle d'évolution}

La boucle tourne 10 fois et suit toujours le même enchaînement :

\begin{enumerate}
  \item \textbf{Génération.}  On envoie un prompt à \texttt{gemini-2.5-flash-lite}
        qui décrit le problème, les meilleures fonctions trouvées jusqu'ici, et les
        conjectures que personne n'a encore réfutées.  Le LLM renvoie une fonction
        Python \texttt{heuristic\_score(G, inv, conjecture)}.
  \item \textbf{Compilation.}  On essaie d'exécuter le code.  S'il contient une
        erreur de syntaxe, on passe à l'itération suivante.
  \item \textbf{Évaluation.}  La fonction guide la recherche sur 25 conjectures
        choisies aléatoirement (60~s max par conjecture, en parallèle).  On mesure
        le coût total (temps de réfutation + 120~s de pénalité par échec).
  \item \textbf{Sélection.}  On garde les 3 meilleures fonctions (coût le plus bas).
  \item \textbf{Feedback.}  Ces 3 fonctions sont incluses dans le prompt suivant,
        avec la liste des conjectures toujours non réfutées.
\end{enumerate}

\subsection{Intégration avec l'heuristique simple}

La Partie~2 réutilise tout le code de la Partie~1 : même génération de population,
mêmes mutations, même réparation des graphes.  On remplace uniquement la fonction
qui évalue si un graphe est ``prometteur''.  La vérification finale (violation~$> 0$)
reste celle de la Partie~1, ce qui garantit que les contre-exemples produits sont
bien réels.

% =========================================================
%  4. RÉSULTATS EXPÉRIMENTAUX
% =========================================================
\section{Résultats expérimentaux}

\subsection{Performance globale}

\begin{table}[h]
  \centering
  \begin{tabular}{lr}
    \toprule
    \textbf{Métrique} & \textbf{Valeur} \\
    \midrule
    Conjectures réfutées & 99 / 100 \\
    Score total (pénalité incluse) & 221,51 s \\
    Temps moyen de réfutation & 1,03 s \\
    Temps minimal & 0,032 s \\
    Temps maximal & 29,21 s \\
    Réfutées en moins d'1 s & 85 / 99 \\
    Réfutées par la passe atlas & 72 / 99 \\
    \bottomrule
  \end{tabular}
  \caption{Résultats de la dernière exécution (11 mai 2026).}
  \label{tab:results}
\end{table}

La passe atlas est de loin la plus efficace : elle réfute 72 conjectures sur 99
avant même que la boucle de recherche ne commence.  Les 27 restantes sont
réfutées par la recherche locale, la plus longue prenant 29~s.  Au total,
85\% des conjectures réfutées l'ont été en moins d'une seconde.

\subsection{La conjecture non réfutée}

La seule conjecture qu'on n'a pas réussi à réfuter est :

\begin{quote}
  \textit{Si $G$ est connexe et sans griffe, $x = \mathrm{total\_domination\_number}$
  et $y = \mathrm{independence\_number}$, alors $y(G) \leq 1 + x(G)$.}
\end{quote}

Le problème est double : la contrainte ``sans griffe'' réduit beaucoup l'espace
des graphes explorables, et le nombre de domination totale est coûteux à calculer.
En pratique, l'inégalité semble toujours vraie pour ces graphes — il est possible
que cette conjecture soit vraie et donc non réfutable.

\subsection{Distribution des temps de réfutation}

Les temps de réfutation sont très déséquilibrés : 85 conjectures en moins d'1~s
(trouvées dans l'atlas), et 14 entre 1 et 30~s (trouvées par la recherche locale).
Ce déséquilibre s'explique simplement : certaines conjectures ont des contre-exemples
très petits ($K_2$, $K_3$), d'autres demandent des graphes plus élaborés.

\subsection{Contribution de FunSearch}

FunSearch n'a pas changé le score final (déjà à 99/100 avec l'heuristique simple).
Sur certaines conjectures difficiles, les fonctions générées par le LLM ont
légèrement accéléré la recherche, mais l'impact global reste limité.

% =========================================================
%  5. DISCUSSION SCIENTIFIQUE
% =========================================================
\section{Discussion scientifique}

\subsection{Quelles conjectures sont faciles à réfuter ?}

Les plus faciles sont celles qui portent sur des invariants liés à la densité :
\texttt{clique\_number}, \texttt{average\_degree}, \texttt{triangle\_number},
\texttt{maximum\_degree}.  Ces invariants montent très vite dans les graphes
denses, donc un simple $K_2$ ou $K_3$ suffit souvent.  Par exemple, la conjecture
$\mathrm{clique\_number}(G) \leq -1/2 + 15/8 \cdot \bar{d}(G)$ est réfutée
par $K_2$ en 0,06~s (violation de 0,625).

Les conjectures sur les invariants de distance (diamètre, proximité) dans la
classe \texttt{connected} sont aussi plus faciles que prévu : les chemins et
cycles, présents dans l'atlas, donnent souvent des contre-exemples immédiats.

\subsection{Quels invariants sont difficiles à manipuler ?}

Trois catégories posent problème :

\begin{itemize}
  \item \textbf{total\_domination\_number avec claw\_free} : l'espace des
        graphes sans griffe est très restreint, et le nombre de domination
        totale change peu d'une mutation à l'autre.  C'est la seule conjecture
        non réfutée.
  \item \textbf{second\_smallest\_laplace\_eigenvalue} (connectivité algébrique) :
        cet invariant est très peu sensible aux petites modifications du graphe,
        ce qui rend la recherche locale lente à converger.
  \item \textbf{largest\_distance\_eigenvalue} : invariant rare, absent de NetworkX,
        qui demande de construire la matrice des distances et de calculer ses
        valeurs propres à chaque évaluation — c'est coûteux.
\end{itemize}

\subsection{Quelles mutations sont efficaces ?}

Les mutations les plus utiles dans notre expérience sont :

\begin{enumerate}
  \item \textbf{Le kick (4 mutations d'un coup)} : indispensable pour sortir
        des blocages.  Sans lui, la recherche reste coincée sur des graphes
        dont le score de violation est légèrement négatif et ne progresse plus.
  \item \textbf{Augmenter le diamètre} (accrocher un nouveau sommet à une feuille) :
        une modification simple qui a un effet immédiat sur le diamètre, la
        proximité et la remoteness à la fois.
  \item \textbf{Les seeds $K_n +$ bras} : ces graphes (une clique avec un chemin
        accroché) combinent haute densité locale et longue portée, ce qui les
        rend utiles pour beaucoup de conjectures en même temps.
\end{enumerate}

Les mutations aléatoires (30\% du temps) sont aussi utiles : elles empêchent
la recherche de se concentrer trop sur une seule structure.

\subsection{FunSearch améliore-t-il réellement l'heuristique ?}

Dans notre cas, FunSearch n'a pas amélioré le score final : l'heuristique
simple était déjà à 99/100.  Sur des conjectures difficiles prises séparément,
les fonctions du LLM ont parfois guidé la recherche un peu plus vite, mais
l'écart reste faible.

FunSearch serait probablement plus utile si l'heuristique de départ était moins
efficace : c'est quand on ne sait pas comment guider la recherche que le LLM
peut apporter quelque chose de nouveau.

\subsection{L'IA a-t-elle produit des idées utiles ?}

En partie.  Parmi les idées originales générées par le LLM :

\begin{itemize}
  \item Récompenser les graphes dont la densité dépasse un seuil adaptatif
        selon la conjecture, ce qui oriente vers des graphes plus denses.
  \item Utiliser le ratio degré max / degré min comme indicateur de
        l'hétérogénéité du graphe.
  \item Pénaliser les graphes réguliers quand la conjecture fait intervenir
        des invariants non linéaires.
\end{itemize}

En revanche, beaucoup de fonctions générées étaient banales (renvoyer directement
le score de violation) ou carrément invalides syntaxiquement.  Le LLM est utile
comme source d'idées, mais pas fiable pour les appliquer correctement sans
vérification.

% =========================================================
%  CONCLUSION
% =========================================================
\section*{Conclusion}

On a réussi à réfuter \textbf{99 conjectures sur 100} avec un score total de
\textbf{221,51~s}.  L'heuristique simple fait l'essentiel du travail : la passe
atlas réfute 72\% des cas en quelques millisecondes, et la recherche locale avec
mutations ciblées gère le reste.  La seule conjecture non réfutée (sans griffe +
domination totale) semble peut-être vraie.

FunSearch est une piste intéressante, mais dans ce contexte il n'a pas apporté
de gain visible.  Pour la suite, il serait plus prometteur de lui demander de
générer directement des familles de graphes plutôt que des fonctions de score,
ou de l'utiliser sur des conjectures où on n'a aucune idée de comment guider la
recherche à la main.

% =========================================================
%  BIBLIOGRAPHIE
% =========================================================
\newpage
\bibliographystyle{plain}
\begin{thebibliography}{9}

\bibitem{romera2024mathematical}
B.~Romera-Paredes, M.~Barekatain, A.~Novikov \emph{et al.}
\newblock Mathematical discoveries from program search with large language models.
\newblock \emph{Nature}, 625:468--475, 2024.

\bibitem{hagberg2008exploring}
A.~Hagberg, P.~Swart, D.~S.~Chult.
\newblock Exploring network structure, dynamics, and function using {NetworkX}.
\newblock In \emph{Proceedings of the 7th Python in Science Conference
  (SciPy2008)}, pages 11--15, 2008.

\bibitem{glover1989tabu}
F.~Glover.
\newblock Tabu search — Part I.
\newblock \emph{ORSA Journal on Computing}, 1(3):190--206, 1989.

\bibitem{bondy1976graph}
J.~A.~Bondy, U.~S.~R.~Murty.
\newblock \emph{Graph Theory with Applications}.
\newblock Macmillan, 1976.

\bibitem{aouchiche2013proximity}
M.~Aouchiche, P.~Hansen.
\newblock Proximity and remoteness in graphs: results and conjectures.
\newblock \emph{Networks}, 58(2):95--102, 2013.

\end{thebibliography}

\end{document}
